\begin{tikzpicture}[
    font=\small,
    >=Stealth,
    gridcell/.style={
        draw=dark!20,   % Leichtes Dark für sichtbare Linien
        fill=gray,      % F3F4F6 ist ohnehin schon ein sehr helles Grau
        line width=0.4pt
    },
    selected/.style={
        draw=amber,
        fill=amber!15,
        line width=0.9pt
    },
    nonselected/.style={
        draw=amber!40,
        fill=gray,
        line width=0.9pt
    },
    testpoint/.style={
        circle,
        fill=teal,
        draw=white,
        line width=0.4pt,
        inner sep=1.6pt
    }
]

    % -------------------------------------------------------------------------
    % Parameter & Koordinaten
    % -------------------------------------------------------------------------
    \def\xmin{0} \def\xmax{7}
    \def\ymin{0} \def\ymax{5}
    \coordinate (xeq) at (3.5, 2.5);

    % -------------------------------------------------------------------------
    % 1. Achsen
    % -------------------------------------------------------------------------
    \draw[->, dark, line width=0.7pt] (\xmin-0.2, \ymin) -- (\xmax+0.5, \ymin) 
        node[right, font=\small\bfseries] {$x_1$};
    \draw[->, dark, line width=0.7pt] (\xmin, \ymin-0.2) -- (\xmin, \ymax+0.5) 
        node[above, font=\small\bfseries] {$x_2$};

    % -------------------------------------------------------------------------
    % 2. Gesamtes Partitionierungsgitter (Hintergrund)
    % -------------------------------------------------------------------------
    \foreach \x in {0,...,6} {
        \foreach \y in {0,...,4} {
            \draw[gridcell] (\x,\y) rectangle ++(1,1);
        }
    }

    % -------------------------------------------------------------------------
    % 3. Selektierte Zellen (passend zur Ellipse)
    % -------------------------------------------------------------------------
    \foreach \x/\y in {
        1/1, 2/1, 3/1, 4/1, 5/1,
        1/2, 2/2, 3/2, 4/2, 5/2,
        1/3, 2/3, 3/3, 4/3, 5/3
    }{
        \draw[selected] (\x,\y) rectangle ++(1,1);
    }

    % -------------------------------------------------------------------------
    % 4. Sublevel-Set V_rho (Als Ellipse)
    % -------------------------------------------------------------------------
    \begin{scope}
        \fill[navy, fill opacity=0.15] (xeq) ellipse [x radius=2.1, y radius=1.3];
        \draw[navy, line width=1.2pt] (xeq) ellipse [x radius=2.1, y radius=1.3];
    \end{scope}

    % Beschriftungen im Plot
    \node[text=navy, font=\small\bfseries] at ($(xeq) + (-0.9, 0.55)$) {$\mathcal{V}_\rho$};
    \node[above left=2pt and 2pt, font=\small\bfseries, text=dark!70] at (\xmax, \ymax) {$\mathcal{X}_V$};

    % -------------------------------------------------------------------------
    % 5. Auswertungspunkte NUR in Beispiel-Zelle (x=1, y=2)
    %    Exakt auf den 2^d Ecken + Zentrum
    % -------------------------------------------------------------------------
    \node[testpoint] (pt_ll)     at (1.0, 2.0) {};
    \node[testpoint] (pt_lr)     at (2.0, 2.0) {};
    \node[testpoint] (pt_ul)     at (1.0, 3.0) {};
    \node[testpoint] (pt_ur)     at (2.0, 3.0) {};
    \node[testpoint] (pt_center) at (1.5, 2.5) {};

    % -------------------------------------------------------------------------
    % 6. Kompakte Legende
    % -------------------------------------------------------------------------
    \begin{scope}[shift={(0.0, -1.35)}]
        % Zeile 1
        \draw[gridcell] (0.0, 0.4) rectangle ++(0.35, 0.25);
        \node[anchor=west, font=\scriptsize, text=dark] at (0.45, 0.5) {Unselected cell};

        \draw[selected] (3.2, 0.4) rectangle ++(0.35, 0.25);
        \node[anchor=west, font=\scriptsize, text=dark] at (3.65, 0.5) {Selected cell ($\exists x_{\mathrm{eval}} \in \mathcal{V}_\rho$)};

        % Zeile 2
        \fill[navy, fill opacity=0.15, draw=navy, line width=0.8pt] (0.0, -0.08) rectangle ++(0.35, 0.25);
        \node[anchor=west, font=\scriptsize, text=dark] at (0.45, -0.0) {Sublevel set $\mathcal{V}_\rho$};

        \node[testpoint] at (3.37, -0.0) {};
        \node[anchor=west, font=\scriptsize, text=dark] at (3.65, -0.0) {Evaluation points ($x_{\mathrm{eval}}$)};
    \end{scope}

\end{tikzpicture}